% AI_LOG.tex --- AI use log for Mini-Project 2 (Local Weather and Climate)
%
% How to use this in your project:
%   1. Standalone: compile this file directly (pdflatex AI_LOG.tex) to get
%      a preview PDF of just the log.
%   2. Inside your report: copy everything from \section{AI Log} down to
%      the end into your main .tex file (e.g. as Appendix E), or instead
%      put \input{AI_LOG_body} around that same content in a separate file
%      and \input{} it from your main document. Either way, drop the
%      preamble below (down to \begin{document}) since your main report
%      already has one.
%
\documentclass[11pt]{article}

\usepackage[utf8]{inputenc}
\usepackage[T1]{fontenc}
\usepackage[margin=1in]{geometry}
\usepackage{hyperref}
\usepackage{enumitem}
\usepackage{parskip}

\title{AI Log}
\date{}

\begin{document}

\section{AI Log --- Mini-Project 2: Local Weather and Climate}

\textbf{Tool used:} Claude Code (Claude Sonnet 5, Anthropic), used interactively in a
chat/CLI session with direct access to the repository files and a Python environment
(pandas, numpy, matplotlib, scipy).

\textbf{Scope of use:} AI wrote all of the Python code in \texttt{Mini-Lab2/}
(\texttt{climate\_common.py}, \texttt{climate\_plots.py}, \texttt{plot\_climate.py}) across
a single ongoing conversation. Every change was produced in response to a specific request
from the team --- usually a plain-language ask, sometimes a screenshot of a rendered chart or
of the assignment PDF's task list --- and the team reviewed the actual rendered PNGs after
every change, which is what drove several rounds of revision on the same chart (most notably
the pressure box plot, entries~\ref{item:pbox1}--\ref{item:pbox3}).

This log is a factual record of what was asked and done. The ``You understand and can
explain your results'' requirement from the assignment is not something the AI can certify
on the team's behalf --- before submitting, the team should be able to explain, without this
log open, what each chart shows and why (see Section~\ref{sec:open-items}).

\subsection*{Log of prompts and actions}

Each entry: what was asked $\rightarrow$ what the AI did $\rightarrow$ what happened to it
(accepted / modified / rejected) and why.

\begin{enumerate}[leftmargin=*, itemsep=0.8em]

\item \textbf{``Can you make a generic plotter, that can plot the different plots from the
folder `Oslo'? I want to be able to change the folder to another place, and it should find
the files in that folder and plot it.''}
$\rightarrow$ Built three modules: \texttt{climate\_common.py} (discovers every CSV in a
place's folder and parses the seklima.met.no export format --- \texttt{;}-delimited, UTF-8
BOM, Norwegian decimal commas, \texttt{-} for missing, a trailing source-note row); a value
column is matched to a physical quantity (temperature/precipitation/pressure/wind) by its
\emph{Norwegian column header}, not the filename, so the same code works unchanged for any
other place's folder. First chart: a shaded climate-normal band for daily max/min temperature
with 2025's actual values overlaid, matching Part 1's two listed tasks (climate profile +
2025-vs-normal comparison) in one figure.
$\rightarrow$ \textbf{Accepted}; the team confirmed by photo that this matched the actual
assignment wording.

\item \textbf{``Can you add a plot that shows temperature something like this as well?''}
(a Carbon Brief ``one line per year'' chart) $\rightarrow$ \textbf{``Can you use the daily
temps instead of finding the monthly?''}
$\rightarrow$ First built a monthly-resolution version, then immediately rebuilt it at daily
resolution once asked: one line per year of 7-day-smoothed daily temperature anomaly against
a day-of-year climatology, coloured by decade, most recent years picked out in bold.
$\rightarrow$ \textbf{Accepted} (daily version kept; this chart was later dropped in the
Part-3/cleanup pass, entry~\ref{item:cleanup}).

\item \textbf{``Seems to be too low resolution to be useful\ldots the bands just show the
max\ldots do you have a suggestion?''} (on the temperature climate-profile chart)
$\rightarrow$ Diagnosed the problem: the shaded band was the all-time record min/max, which
is so wide it almost never excludes anything. Replaced it with a smoothed 10th--90th
percentile ``typical range'' band, with the true record extremes kept as thin dotted
reference lines.
$\rightarrow$ \textbf{Accepted.}

\item \textbf{``Can you color the section `january' gray, `february' white, and so on\ldots
on days 1--365?''}
$\rightarrow$ Added a shared \texttt{\_shade\_months()} helper (alternating light-gray/white
vertical bands keyed to calendar months) and applied it to every day-of-year chart.
$\rightarrow$ \textbf{Accepted.}

\item \textbf{``Does not say a lot?''} (on the precipitation monthly bar chart, mean $\pm$ 1
std)
$\rightarrow$ Replaced the bar-plus-error-bar chart with a box plot (10th--90th percentile
whiskers) with every individual year's monthly total scattered on top --- mean $\pm$ std
hides that precipitation is right-skewed and can't go below zero.
$\rightarrow$ \textbf{Accepted.}

\item \textbf{``Can you try to fit a gamma distribution to the plots? Maybe create new plots
so they don't pollute the already existing plots.''}
$\rightarrow$ Added two new charts (existing ones untouched): a daily wet-day rainfall-amount
histogram with a fitted \texttt{scipy.stats.gamma} curve plus a Q--Q plot, and a 12-panel grid
fitting gamma separately to each calendar month's year-to-year totals. While building the
second one, found and fixed a real bug: \texttt{pandas}' \texttt{.sum()} silently returns 0
(not NaN) for an all-missing group, which had been feeding a spurious 0~mm month into the
fit; separately, one genuinely dry month (April 1974, confirmed against the raw CSV) had to
be excluded from its own fit since a continuous distribution assigns exactly 0 zero
probability.
$\rightarrow$ \textbf{Accepted.}

\item \textbf{``Can you name things? As in mark 2025/26, and name what the box-plot
shows?''}
$\rightarrow$ Added 2025/2026 highlight markers (hollow if the month was still in progress)
and explicit legend labels for every box element. While doing so, found the existing legend
mislabeled the box as ``10th--90th percentile'': with \texttt{whis=(10, 90)}, matplotlib
actually draws the box as the 25th--75th percentile (IQR) and only the \emph{whiskers} at the
10th/90th --- fixed on both the precipitation and (for consistency) the pressure box plot.
$\rightarrow$ \textbf{Accepted.}

\item \textbf{``Maybe color code it as well? Adding a gradient?''} \textbf{``Can you add an
average as well to the box plot?''}
$\rightarrow$ Added a year-coloured (coolwarm) scatter with a colourbar to the precipitation
box plot, and a mean marker (\texttt{showmeans}) to both the precipitation and pressure box
plots, to make the skew (mean pulled above the median in wet months) visible directly.
$\rightarrow$ \textbf{Accepted.}

\item \textbf{``Can you add a variable where we can choose the first year, as in a lower
bound?''}
$\rightarrow$ Added \texttt{START\_YEAR} to \texttt{plot\_climate.py}; every chart is rebuilt
from only that year onward when set (\texttt{None} keeps the full record). The team later set
this to \texttt{1991} themselves.
$\rightarrow$ \textbf{Accepted.}

\item \textbf{``Can you change the pressure whiskers to max and min in the reference period,
instead of the 1.5 IQR? Can you also make the legends describe what they show, as in `25--75
percentile (reference years)'?''}
$\rightarrow$ Restricted the box/whisker statistic on both the precipitation and pressure box
plots to the 1991--2020 reference period specifically (the scatter/highlight layer kept using
the full record), and switched the pressure whiskers from the default 1.5$\times$IQR (Tukey)
rule to the true reference-period min/max, so anything outside the whiskers is genuinely
unprecedented rather than a statistical-outlier cutoff. Every legend entry across both charts
was relabelled with its exact period, e.g. ``25th--75th percentile (1991--2020)''.
$\rightarrow$ \textbf{Accepted.}

\item \textbf{``I need a graph to see if the precipitation has increased over time. Maybe a
spaghetti plot? Cumulative plot?''}
$\rightarrow$ Built a two-panel chart: annual totals with a 10-year moving average on top, and
the cumulative departure from the long-term mean (a ``mass curve'', standard hydrological
technique) on the bottom --- a single straight trend line through noisy annual bars is easy
to misread as ``no trend'' even when a real one exists.
$\rightarrow$ \textbf{Accepted}, later simplified (entry~\ref{item:masscurve}).

\item \textbf{``Can you add a poly fit and poly val to the graph? Can you also try to create
a spaghetti plot for rain as well?''}
$\rightarrow$ Added a degree-2 polynomial fit (\texttt{np.polyfit}/\texttt{np.polyval})
alongside the moving average on the precipitation trend chart; and, as a new chart, one line
per year of \emph{cumulative} precipitation through the year (day-of-year on the x-axis,
coloured by decade) --- raw daily rainfall is too spiky for a year-over-year spaghetti plot,
but its running total is smooth and monotonic.
$\rightarrow$ \textbf{Accepted.}

\item \textbf{``Can you do a suitable test to compare the distributions of the climate
profile and the precipitation of 2025?''}
$\rightarrow$ Built \texttt{climate\_stats.py}: a two-sample Kolmogorov--Smirnov test
(primary) plus Mann-Whitney U (secondary), explicitly \emph{not} a t-test, since daily
precipitation is right-skewed and a t-test's normality assumption/means-only comparison would
be the wrong tool; plus an ECDF chart showing what the KS statistic means visually.
\textbf{Judgement call, stated to the team at the time:} the request was ambiguous between
temperature and precipitation --- read as precipitation, given the conversation's focus.
$\rightarrow$ \textbf{Accepted at the time, later removed entirely}
(entry~\ref{item:cleanup}) as not on the assignment's task list.

\item \textbf{``Can you make the plot more pretty? Can you add the 2025 and 2026 markers to
Oslo\_pressure\_monthly\_boxplot as well?''}
$\rightarrow$ Polished the ECDF chart (shaded gap between curves, an annotated arrow at the
KS statistic instead of an all-but-invisible line, percentage axis); added 2025/2026 markers
(monthly-mean based, first version) to the pressure box plot.
$\rightarrow$ \textbf{Accepted at the time} (ECDF chart later removed; the pressure-box
markers went through two more rounds, entries~\ref{item:pbox1}--\ref{item:pbox3}).

\item \textbf{``I feel like the pressure graphs don't really answer these criteria?''}
(screenshot of Part 3a's actual task text: ``explore variations during a month'' /
``investigate whether pressure has changed over time'' / the framing question ``is this
pressure normal for X?'')
$\rightarrow$ Pointed out that the two existing pressure charts answered the two listed
tasks but not the section's own framing question. Added \texttt{plot\_pressure\_vs\_normal},
mirroring the Part 1 temperature chart exactly: a smoothed typical-range band with the actual
year's daily pressure overlaid.
$\rightarrow$ \textbf{Accepted.}

\item \textbf{``Can you create a plot of the total rainfall with a subplot for every month,
to see if there is an increase in that specific month over time?''} \textbf{``Can you add a
poly fit and value?''}
$\rightarrow$ Built a 12-panel grid (one panel per calendar month) with each month's own
linear trend (coloured red/blue by direction) and a degree-2 polynomial fit. This revealed
that the increase is not uniform across the year --- July/August/September/December trend
clearly wetter, while April and November trend \emph{drier}. (Fixed one rendering bug along
the way: an out-of-bounds \texttt{fig.legend()} was being silently clipped off the saved PNG;
replaced with the existing caption-text convention.)
$\rightarrow$ \textbf{Accepted.}

\item \label{item:cleanup} \textbf{``Here are the actual tasks. Can you check what we have,
improve it, and remove the plots not needed?''} (assignment PDF attached)
$\rightarrow$ Cross-checked every chart against the task list. Two decisions could only be
made by the team, so these were asked directly rather than assumed:
\begin{itemize}[leftmargin=*]
\item \emph{Part 3 says choose Pressure \textbf{or} Wind} --- both were being generated.
Team chose \textbf{Pressure}.
\item \emph{Which ``bonus'' charts (not on the task list) to keep} --- of the temperature
spaghetti, precipitation cumulative spaghetti, the KS-test/ECDF chart, and the older
\texttt{precip\_annual\_trend} (superseded by the cumulative-trend chart), the team kept only
the \textbf{precipitation cumulative spaghetti}.
\end{itemize}
Removed \texttt{plot\_wind\_monthly\_climatology},
\texttt{plot\_temperature\_daily\_spaghetti},
\texttt{plot\_precip\_ecdf\_2025\_vs\_normal} together with the whole
\texttt{climate\_stats.py} module, and \texttt{plot\_precip\_annual\_trend}, plus their stale
PNGs. Figure count went from 13 to 9.
$\rightarrow$ \textbf{Accepted.}

\item \label{item:masscurve} \textbf{``Can you remove this part of the rain plot?''}
(screenshot of the cumulative-anomaly ``mass curve'' bottom panel)
$\rightarrow$ Removed that panel; the precipitation trend chart is a single panel again
(annual bars + moving average + polynomial fit + long-term mean line).
$\rightarrow$ \textbf{Accepted.}

\item \label{item:pbox1} \textbf{``\texttt{Oslo\_pressure\_annual\_trend} doesn't really tell
us something about `has it changed over time'. Also,
\texttt{Oslo\_pressure\_monthly\_boxplot} has whiskers with no datapoint at the ends?''}
$\rightarrow$ Rebuilt the pressure annual-trend chart the same way the precipitation one had
already been fixed (moving average + polynomial fit through the noisy annual points). For the
box plot: diagnosed the ``floating whisker'' as a resolution mismatch --- the whiskers were
built from \emph{daily} values while the highlighted-year dots were \emph{monthly means},
which can mathematically never reach a daily extreme. First fix attempt: plot each
highlighted year's actual daily values instead of one monthly-mean marker. (Also caught and
fixed an unrelated stray title bug, ``pressure variation within a \emph{year}'', that had crept
in during an earlier edit and should have read ``month''.)
$\rightarrow$ \textbf{Modified next round} (entry~\ref{item:pbox2}).

\item \label{item:pbox2} \textbf{``I actually liked it more before, now it seems too
crowded? And there are still no datapoints at the max/min.''}
$\rightarrow$ Realised the daily-value swarm was both busier than the earlier version \emph{and}
still didn't fully fix the mismatch (a highlighted year's own extreme rarely coincides with
the 30-year record extreme). Tried removing the per-year overlay entirely, back to a plain
box plot.
$\rightarrow$ \textbf{Rejected} --- the team wanted the coloured scatter kept, just without
the mismatch (see next entry).

\item \label{item:pbox3} \textbf{``I want the scatter to a degree, like it was in the
attached image, but there are no datapoints at the ends of the whiskers.''}
$\rightarrow$ Root-caused properly this time: rebuilt the box \emph{and} whiskers from the
same array being scattered (each year's mean pressure for that month), instead of from daily
values, so the whisker tips are always literally the most extreme plotted dot. Restored the
year-colour gradient and colourbar.
$\rightarrow$ \textbf{Accepted.} Note: this changes what the chart technically measures ---
see Section~\ref{sec:open-items}.

\item \textbf{``A lot of the legends here overlap with the graph\ldots maybe add a little to
the top of the y-lim?''} (4 named charts) \textbf{``The first two could need another nudge
upwards?''}
$\rightarrow$ Added a shared \texttt{\_add\_top\_margin()} helper and applied it with a margin
tuned per chart (temperature: a small 6\% nudge; the three pressure charts: 12--20\%, the two
noisiest ones increased further on request).
$\rightarrow$ \textbf{Accepted.}

\item \textbf{``Can you condense all the things I have asked for here\ldots create a tex
file with that summary?''}
$\rightarrow$ This file.

\end{enumerate}

\subsection*{Open items for the team}
\label{sec:open-items}

A few places where the AI made a call, or where a chart's exact meaning is easy to
misstate, that the team should actively confirm and be ready to explain --- not just inherit
silently:

\begin{itemize}[leftmargin=*]
\item \textbf{\texttt{Oslo\_pressure\_monthly\_boxplot} shows year-to-year spread of each
month's \emph{mean} pressure, not day-to-day spread within a month.} That day-to-day
question is answered instead by \texttt{Oslo\_pressure\_vs\_normal} (actual daily pressure
against a smoothed typical-range band). Be ready to explain why the box plot uses one year's
mean per point rather than individual days --- it's because a box built from daily values
can never have its whiskers coincide with a scattered point at monthly-mean resolution (see
entries~\ref{item:pbox1}--\ref{item:pbox3}).
\item \textbf{\texttt{START\_YEAR = 1991}} restricts every current chart to 1991 onward, even
though the raw temperature/precipitation records go back to 1937 and pressure/wind to 1957.
Confirm this is the intended window for the report, not an oversight.
\item \textbf{Part 3 uses Pressure, not Wind} (team's explicit choice, entry~\ref{item:cleanup}).
The wind columns are still in the raw \texttt{wind\_pressure\_Oslo.csv} if the team changes
its mind, but the plotting function for it was deleted (retrievable from git history).
\item \textbf{The gamma-fit ``amount of precipitation'' teamwork question} is answered in
code by \emph{two} different definitions --- daily amount on wet days
(\texttt{Oslo\_precip\_daily\_gamma\_fit}) and monthly total
(\texttt{Oslo\_precip\_monthly\_gamma\_fit}) --- but the report still needs a written sentence
or two making that comparison explicit; the code doesn't write the discussion for you.
\item Appendix D (AI Declaration Form) and Appendix F (team-efficiency reflections) are not
things the AI can produce on the team's behalf --- still to do.
\end{itemize}

\end{document}
